\documentclass[a4paper]{article}

\usepackage{graphicx}
\usepackage{amsmath,amssymb}
\usepackage{minted}
\usepackage{multirow}
\usepackage{float}
\usepackage{todonotes}
\usepackage{forest}
\usepackage{xstring}
\usepackage{xcolor}
\usepackage[most]{tcolorbox}
\usepackage{xparse}
\usepackage{natbib}

\usetikzlibrary{graphs,graphdrawing}
\usegdlibrary{layered}

% for external graphviz
\input{/home/donkarlo/Dropbox/repo/nd_language_ai_project/src/nd_language_ai/formal/computer/programming/tex/latex/code_clips/external_graphviz.tex}

% for tags
\input{/home/donkarlo/Dropbox/repo/nd_language_ai_project/src/nd_language_ai/formal/computer/programming/tex/latex/parts/control_sequence/kind/semtag/semtag.tex}

% for links
\input{/home/donkarlo/Dropbox/repo/nd_language_ai_project/src/nd_language_ai/formal/computer/programming/tex/latex/code_clips/seehere.tex}

\title{Neural coding}
\date{}
\author{Mohammad Rahmani}

\begin{document}
    \maketitle
    Neural coding (or neural representation) refers to the relationship between a stimulus and its respective neuronal responses, and the signaling relationships among networks of neurons in an ensemble. \seeherefooter{https://en.wikipedia.org/wiki/Neural\_coding}
    
    \cite{shimazaki-2025-neural-coding-foundational-concepts-statistical-formulations-and-recent-advances} explains neural coding as the process by which stimuli are transformed into neural activity that supports perception and behavior, combining classical concepts with statistical encoding and decoding frameworks.
    \bibliographystyle{plainnat}
    \bibliography{/home/donkarlo/Dropbox/repo/nd_shared_research/bibliography/bibliography.bib}
\end{document}